\documentclass[11pt,a4paper]{ctexart}

\usepackage[margin=1in]{geometry}
\usepackage{amsmath,amssymb,amsthm,mathtools,bm}
\usepackage{hyperref}
\usepackage{enumitem}
\usepackage{booktabs}

\hypersetup{
  colorlinks=true,
  linkcolor=blue,
  citecolor=blue,
  urlcolor=blue
}

% ----------------------------
% Macros
% ----------------------------
\newcommand{\RR}{\mathbb{R}}
\newcommand{\ZZ}{\mathbb{Z}}
\newcommand{\CC}{\mathbb{C}}

\newcommand{\Enc}{\mathsf{Enc}}
\newcommand{\Auto}[1]{\mathsf{Auto}_{#1}}
\newcommand{\Add}{\mathsf{Add}}
\newcommand{\Sub}{\mathsf{Sub}}
\newcommand{\MulPt}[1]{\mathsf{MulPt}_{#1}}
\newcommand{\Round}{\operatorname{Round}}

\newcommand{\ct}{\mathbf{ct}}
\newcommand{\cte}{\mathbf{ct}_{\mathrm{even}}^{\uparrow}}
\newcommand{\ctoy}{\mathbf{ct}_{\mathrm{odd},Y}^{\uparrow}}
\newcommand{\cto}{\mathbf{ct}_{\mathrm{odd}}^{\uparrow}}

\newcommand{\Rbig}{\mathcal{R}_N^{\RR}}
\newcommand{\Rint}{\mathcal{R}_N^{\ZZ}}
\newcommand{\Ssmall}{\mathcal{S}_n}

\newcommand{\rep}{\mathrel{\vDash_{\Delta}}}

% ----------------------------
% Theorem environments
% ----------------------------
\newtheorem{definition}{定义}[section]
\newtheorem{proposition}[definition]{命题}
\newtheorem{theorem}[definition]{定理}
\newtheorem{lemma}[definition]{引理}
\newtheorem{corollary}[definition]{推论}
\newtheorem{remark}[definition]{注}

% ----------------------------
% Title
% ----------------------------
\title{在 CKKS 密文上实现多项式的 FFT 式奇偶拆分与重构\\
\large 一份自洽且可直接落地的推导说明}
\author{}
\date{}

\begin{document}

\maketitle


\tableofcontents
\newpage


\begin{abstract}
本文整理并严格推导一个可直接在 CKKS 密文上执行的“FFT 式奇偶拆分 / 反向重构”过程。设
\[
P(Y)=\sum_{j=0}^{2n-1}p_jY^j \in \RR[Y]/(Y^{2n}+1),
\]
我们希望在不经过 CoeffToSlot / SlotToCoeff 的前提下，直接从一个加密 \(P(Y)\) 的 CKKS 密文出发，得到两个分别对应
\[
P_{\mathrm{even}}(T)=\sum_{i=0}^{n-1}p_{2i}T^i,
\qquad
P_{\mathrm{odd}}(T)=\sum_{i=0}^{n-1}p_{2i+1}T^i
\]
的密文；并进一步给出从这两个分支密文重构原始密文的完整公式。全文采用多项式视角而非槽向量视角，因为这样可以直接使用 CKKS 的标准操作：自同构、加法、减法与明文--密文乘法。文中所有代数等式都在明文多项式层面精确成立；对应的密文操作在 CKKS 语义下可直接实现。
\end{abstract}

\section{目标与方法选择}

本文的目标是严格说明以下两件事：

\begin{enumerate}[label=\arabic*.]
    \item \textbf{正向拆分：} 从一个加密
    \[
    P(Y)=\sum_{j=0}^{2n-1}p_jY^j
    \]
    的 CKKS 密文出发，同态地得到两个分支密文，分别对应 FFT 递归中的偶部分和奇部分；
    
    \item \textbf{反向重构：} 从这两个分支密文出发，同态地恢复原始的 \(P(Y)\) 密文。
\end{enumerate}

需要强调的是，这里我们\textbf{不采用}直接把系数向量
\[
(p_0,p_1,\dots,p_{2n-1})^\top
\]
看成一个长度 \(2n\) 向量，再用矩阵 \(E,O\) 去做抽取的方案。原因是：标准 CKKS 的“矩阵--密文乘法”作用对象本质上是\textbf{槽向量表示}，而这里我们要处理的是\textbf{多项式系数表示}。如果走矩阵 \(E,O\) 的路线，通常需要先做一次 CoeffToSlot，把系数信息变成槽向量；那样虽然理论上也可行，但不如直接在多项式域里做来得自然。

因此，本文选择如下更适合 CKKS 密文域的路线：
\[
P(Y)=P_{\mathrm{even}}(Y^2)+Y\,P_{\mathrm{odd}}(Y^2),
\]
并结合
\[
P(-Y)=P_{\mathrm{even}}(Y^2)-Y\,P_{\mathrm{odd}}(Y^2)
\]
来实现分裂与重构。

\section{记号、环与 CKKS 语义约定}

\subsection{基本参数与环}

固定一个 2 的幂 \(N\ge 1\)，以及一个满足 \(n\mid N\) 且 \(n < N\) 的 2 的幂。定义
\[
k=\frac{N}{2n}\in \ZZ_{>0}.
\]

考虑两个环：

\begin{align*}
\Ssmall &= \RR[Y]/(Y^{2n}+1),\\
\Rbig &= \RR[X]/(X^N+1).
\end{align*}

通过代换
\[
Y=X^k
\]
把 \(\Ssmall\) 嵌入到 \(\Rbig\) 中。由于
\[
Y^{2n}=(X^k)^{2n}=X^{2nk}=X^N=-1,
\]
所以该嵌入是良定义的。

在 CKKS 中，实际明文多项式（经过缩放与取整后）属于整数环
\[
\Rint=\ZZ[X]/(X^N+1).
\]
但本文中的代数推导主要发生在实系数环 \(\Rbig\) 中；最后再通过 CKKS 编码把实系数对象送入 \(\Rint\)。

\subsection{多项式与偶奇分支}

设
\[
P(Y)=\sum_{j=0}^{2n-1}p_jY^j \in \Ssmall.
\]

定义两个多项式：
\begin{align}
P_{\mathrm{even}}(T) &= \sum_{i=0}^{n-1}p_{2i}T^i, \label{eq:def-even}\\
P_{\mathrm{odd}}(T) &= \sum_{i=0}^{n-1}p_{2i+1}T^i. \label{eq:def-odd}
\end{align}

于是
\[
P(Y)=P_{\mathrm{even}}(Y^2)+Y\,P_{\mathrm{odd}}(Y^2).
\]

这里
\[
P_{\mathrm{even}}(T)
\]
的系数向量正是
\[
(p_0,p_2,\dots,p_{2n-2}),
\]
而
\[
P_{\mathrm{odd}}(T)
\]
的系数向量正是
\[
(p_1,p_3,\dots,p_{2n-1}).
\]

因此，这两个多项式正对应 FFT 递归中“偶数项系数向量”和“奇数项系数向量”。

\subsection{CKKS 语义记号}

为了把推导写得清楚，同时不被实现细节淹没，我们约定如下的\textbf{语义记号}。

\begin{definition}[密文表示的消息]
设当前 CKKS scale 为 \(\Delta\)。若一个密文 \(\ct\) 解密后得到的明文多项式为
\[
m(X)=\Round\!\bigl(\Delta F(X)\bigr)+e(X)\in \Rint,
\]
其中 \(F(X)\in \Rbig\)，而 \(e(X)\) 是 CKKS 中通常的小噪声项，则记为
\[
\ct \rep F(X).
\]
这表示：\(\ct\) 在当前 scale 下\textbf{表示}消息 \(F(X)\)。
\end{definition}

这里 \(\Round\) 表示逐系数取最近整数。于是：

\begin{itemize}
    \item 当写 \(\ct \rep F(X)\) 时，真正被加密的是 \(\Round(\Delta F(X))\)（再加上通常的小加密噪声）；
    \item 本文中的代数等式都作用在 \(F(X)\) 这个\textbf{被表示的消息}层面；
    \item 若实现中需要做 rescale、modulus switching 或 scale matching，则默认按标准 CKKS 方式处理，使得\(\ct\) 所表示的消息不变。为集中讨论代数正确性，这些实现细节在公式中省略。
\end{itemize}

\subsection{本文使用的同态操作}

我们只使用以下四类标准操作。

\begin{enumerate}[label=(\roman*)]
    \item \textbf{自同构：}
    \[
    \Auto{a}(\ct), \qquad a \text{ 为奇数},
    \]
    表示对密文应用环自同构
    \[
    \sigma_a:\Rbig\to\Rbig,\qquad f(X)\mapsto f(X^a).
    \]

    \item \textbf{加法：}
    \[
    \Add(\ct_1,\ct_2).
    \]

    \item \textbf{减法：}
    \[
    \Sub(\ct_1,\ct_2).
    \]

    \item \textbf{明文--密文乘法：}
    \[
    \MulPt{u(X)}(\ct),
    \]
    其中 \(u(X)\in \Rbig\) 是一个明文多项式（例如常数 \(1/2\)、单项式 \(X^k\) 或 \(-X^{N-k}\)）。
\end{enumerate}

这些操作在消息层面满足以下标准语义：

\begin{proposition}[操作的消息语义]
设 \(\ct,\ct_1,\ct_2\) 为 CKKS 密文，\(u(X)\in\Rbig\)，\(a\) 为奇数。
\begin{enumerate}[label=(\arabic*)]
    \item 若 \(\ct \rep F(X)\)，则
    \[
    \Auto{a}(\ct)\rep F(X^a).
    \]
    \item 若 \(\ct_1\rep F_1(X)\) 且 \(\ct_2\rep F_2(X)\)，则
    \[
    \Add(\ct_1,\ct_2)\rep F_1(X)+F_2(X),
    \]
    \[
    \Sub(\ct_1,\ct_2)\rep F_1(X)-F_2(X).
    \]
    \item 若 \(\ct\rep F(X)\)，则
    \[
    \MulPt{u(X)}(\ct)\rep u(X)\,F(X).
    \]
\end{enumerate}
\end{proposition}

\begin{remark}
上面的第三条在实现里可能伴随 scale 调整与 rescale；但在“被表示的消息”这一层面，其作用就是乘上 \(u(X)\)。
\end{remark}

\subsection{上标 \(\uparrow\) 的含义}

稍后我们会得到形如
\[
P_{\mathrm{even}}(X^{2k}), \qquad P_{\mathrm{odd}}(X^{2k})
\]
这样的对象。它们虽然仍属于大环 \(\Rbig\)，但只包含 \(X^{2k}\) 的幂，因此实际上落在更小的子环里。为提醒读者“它仍存放在大环中，但语义上来自更小子环”，我们记对应密文为
\[
\cte,\qquad \ctoy,\qquad \cto,
\]
其中上标 \({}^{\uparrow}\) 表示“提升回大环中的表示”。

\section{精确的代数恒等式}

本节所有等式都发生在 \(\Ssmall=\RR[Y]/(Y^{2n}+1)\) 中，是\textbf{精确等式}。

\begin{proposition}[FFT 式偶奇分解]
设
\[
P(Y)=\sum_{j=0}^{2n-1}p_jY^j.
\]
定义 \(P_{\mathrm{even}}, P_{\mathrm{odd}}\) 如 \eqref{eq:def-even} 与 \eqref{eq:def-odd}。则
\begin{equation}
P(Y)=P_{\mathrm{even}}(Y^2)+Y\,P_{\mathrm{odd}}(Y^2). \label{eq:fft-split}
\end{equation}
\end{proposition}

\begin{proof}
直接展开右边：
\begin{align*}
P_{\mathrm{even}}(Y^2)+Y\,P_{\mathrm{odd}}(Y^2)
&=
\sum_{i=0}^{n-1}p_{2i}(Y^2)^i
+
Y\sum_{i=0}^{n-1}p_{2i+1}(Y^2)^i\\
&=
\sum_{i=0}^{n-1}p_{2i}Y^{2i}
+
\sum_{i=0}^{n-1}p_{2i+1}Y^{2i+1}\\
&=
\sum_{j=0}^{2n-1}p_jY^j\\
&=P(Y).
\end{align*}
\end{proof}

\begin{proposition}[以 \(P(-Y)\) 表示偶奇分支]
在 \eqref{eq:fft-split} 的条件下，有
\begin{equation}
P(-Y)=P_{\mathrm{even}}(Y^2)-Y\,P_{\mathrm{odd}}(Y^2). \label{eq:p-minus-y}
\end{equation}
从而
\begin{equation}
P_{\mathrm{even}}(Y^2)=\frac{P(Y)+P(-Y)}{2}, \label{eq:even-half-sum}
\end{equation}
\begin{equation}
Y\,P_{\mathrm{odd}}(Y^2)=\frac{P(Y)-P(-Y)}{2}. \label{eq:odd-half-diff}
\end{equation}
\end{proposition}

\begin{proof}
将 \eqref{eq:fft-split} 中的 \(Y\) 替换为 \(-Y\) 即得 \eqref{eq:p-minus-y}。再把 \eqref{eq:fft-split} 与 \eqref{eq:p-minus-y} 相加、相减即可得到 \eqref{eq:even-half-sum} 与 \eqref{eq:odd-half-diff}。
\end{proof}

\section{正向方向：从 \texorpdfstring{\(\ct_P\)}{ctP} 到两个分支密文}

\subsection{输入}

设输入密文为
\[
\ct_P \rep P(X^k),
\]
其中
\[
P(X^k)=\sum_{j=0}^{2n-1}p_jX^{jk}\in \Rbig.
\]

也就是说，\(\ct_P\) 是一个直接在 CKKS 大环中表示 \(P(Y)\)（经 \(Y=X^k\) 嵌入后）的密文。

\subsection{第一步：构造 \texorpdfstring{\(P(-Y)\)}{P(-Y)} 的密文}

\begin{proposition}[实现 \(Y\mapsto -Y\) 的自同构]
取
\[
a=1+2n.
\]
则 \(a\) 为奇数，因此对应一个合法的环自同构
\[
\sigma_a:\Rbig\to\Rbig,\qquad f(X)\mapsto f(X^a).
\]
并且
\begin{equation}
\sigma_a(X^k)=-X^k. \label{eq:auto-sends-Y-to-minus-Y}
\end{equation}
\end{proposition}

\begin{proof}
由于 \(2n\) 为偶数，所以 \(a=1+2n\) 为奇数。再由 \(k=N/(2n)\) 可得
\[
2nk=N.
\]
因此
\begin{align*}
\sigma_a(X^k)
&=X^{ka}\\
&=X^{k(1+2n)}\\
&=X^{k+2nk}\\
&=X^{k+N}\\
&=X^k\cdot X^N\\
&=X^k\cdot(-1)\\
&=-X^k.
\end{align*}
这就证明了 \eqref{eq:auto-sends-Y-to-minus-Y}。
\end{proof}

据此定义
\begin{equation}
\ct_{-}:=\Auto{1+2n}(\ct_P). \label{eq:def-ct-minus}
\end{equation}

\begin{corollary}
有
\[
\ct_{-}\rep P(-X^k).
\]
\end{corollary}

\begin{proof}
由 \(\ct_P\rep P(X^k)\) 以及命题 \eqref{eq:auto-sends-Y-to-minus-Y}，再用自同构的消息语义即可：
\[
\ct_{-}
=
\Auto{1+2n}(\ct_P)
\rep
P\bigl(\sigma_{1+2n}(X^k)\bigr)
=
P(-X^k).
\]
\end{proof}

\subsection{第二步：得到偶分支}

定义
\begin{equation}
\cte
:=
\MulPt{\frac12}\!\Bigl(\Add(\ct_P,\ct_{-})\Bigr). \label{eq:def-ct-even}
\end{equation}

\begin{theorem}[偶分支的正确性]
由 \eqref{eq:def-ct-even} 得到的密文满足
\[
\cte \rep P_{\mathrm{even}}(X^{2k}).
\]
更具体地，
\[
P_{\mathrm{even}}(X^{2k})=\sum_{i=0}^{n-1}p_{2i}X^{2ik}.
\]
\end{theorem}

\begin{proof}
由 \(\ct_P\rep P(X^k)\) 与 \(\ct_{-}\rep P(-X^k)\)，先得
\[
\Add(\ct_P,\ct_{-})\rep P(X^k)+P(-X^k).
\]
再乘上 \(1/2\)，得到
\[
\cte \rep \frac{P(X^k)+P(-X^k)}{2}.
\]
根据 \eqref{eq:even-half-sum}，把其中的 \(Y\) 替换成 \(X^k\)，即有
\[
\frac{P(X^k)+P(-X^k)}{2}=P_{\mathrm{even}}((X^k)^2)=P_{\mathrm{even}}(X^{2k}).
\]
于是
\[
\cte \rep P_{\mathrm{even}}(X^{2k}).
\]
最后由 \(P_{\mathrm{even}}(T)=\sum_{i=0}^{n-1}p_{2i}T^i\)，代入 \(T=X^{2k}\) 得
\[
P_{\mathrm{even}}(X^{2k})=\sum_{i=0}^{n-1}p_{2i}X^{2ik}.
\]
\end{proof}

\subsection{第三步：得到带 \texorpdfstring{\(Y\)}{Y} 因子的奇分支}

定义
\begin{equation}
\ctoy
:=
\MulPt{\frac12}\!\Bigl(\Sub(\ct_P,\ct_{-})\Bigr). \label{eq:def-ct-oddY}
\end{equation}

\begin{theorem}[带 \(Y\) 因子的奇分支的正确性]
由 \eqref{eq:def-ct-oddY} 得到的密文满足
\[
\ctoy \rep X^k\,P_{\mathrm{odd}}(X^{2k}).
\]
更具体地，
\[
X^k\,P_{\mathrm{odd}}(X^{2k})
=
\sum_{i=0}^{n-1}p_{2i+1}X^{(2i+1)k}.
\]
\end{theorem}

\begin{proof}
由 \(\ct_P\rep P(X^k)\) 与 \(\ct_{-}\rep P(-X^k)\)，先得
\[
\Sub(\ct_P,\ct_{-})\rep P(X^k)-P(-X^k).
\]
再乘上 \(1/2\)，得到
\[
\ctoy \rep \frac{P(X^k)-P(-X^k)}{2}.
\]
根据 \eqref{eq:odd-half-diff}，把其中的 \(Y\) 替换成 \(X^k\)，得到
\[
\frac{P(X^k)-P(-X^k)}{2}
=
X^k\,P_{\mathrm{odd}}((X^k)^2)
=
X^k\,P_{\mathrm{odd}}(X^{2k}).
\]
于是
\[
\ctoy \rep X^k\,P_{\mathrm{odd}}(X^{2k}).
\]
再展开 \(P_{\mathrm{odd}}(T)=\sum_{i=0}^{n-1}p_{2i+1}T^i\)，可得
\[
X^k\,P_{\mathrm{odd}}(X^{2k})
=
X^k\sum_{i=0}^{n-1}p_{2i+1}(X^{2k})^i
=
\sum_{i=0}^{n-1}p_{2i+1}X^{(2i+1)k}.
\]
\end{proof}

\subsection{第四步（可选）：去掉奇分支前面的 \texorpdfstring{\(Y=X^k\)}{Y=Xk}}

有时为了继续递归拆分，我们希望奇分支也写成与偶分支完全同形的形式，即直接得到
\[
P_{\mathrm{odd}}(X^{2k})
\]
而不是
\[
X^k\,P_{\mathrm{odd}}(X^{2k}).
\]

\begin{lemma}[求 \(Y=X^k\) 的逆元]
在 \(\Rbig=\RR[X]/(X^N+1)\) 中，
\begin{equation}
(X^k)^{-1}=-X^{N-k}. \label{eq:Y-inverse}
\end{equation}
\end{lemma}

\begin{proof}
直接计算：
\[
X^k\cdot(-X^{N-k})=-X^N=1.
\]
故
\[
(X^k)^{-1}=-X^{N-k}.
\]
\end{proof}

定义
\begin{equation}
\cto
:=
\MulPt{-X^{N-k}}(\ctoy). \label{eq:def-ct-odd}
\end{equation}

\begin{theorem}[去掉 \(Y\) 因子的奇分支的正确性]
由 \eqref{eq:def-ct-odd} 得到的密文满足
\[
\cto \rep P_{\mathrm{odd}}(X^{2k}).
\]
更具体地，
\[
P_{\mathrm{odd}}(X^{2k})=\sum_{i=0}^{n-1}p_{2i+1}X^{2ik}.
\]
\end{theorem}

\begin{proof}
由上一节可知
\[
\ctoy \rep X^k\,P_{\mathrm{odd}}(X^{2k}).
\]
再乘上 \(-X^{N-k}\)，由 \eqref{eq:Y-inverse} 可得
\[
\cto
=
\MulPt{-X^{N-k}}(\ctoy)
\rep
(-X^{N-k})\cdot X^k\,P_{\mathrm{odd}}(X^{2k})
=
P_{\mathrm{odd}}(X^{2k}).
\]
最后把
\[
P_{\mathrm{odd}}(T)=\sum_{i=0}^{n-1}p_{2i+1}T^i
\]
中的 \(T\) 替换成 \(X^{2k}\) 即得
\[
P_{\mathrm{odd}}(X^{2k})=\sum_{i=0}^{n-1}p_{2i+1}X^{2ik}.
\]
\end{proof}

\subsection{正向方向的最终结论}

\begin{theorem}[正向奇偶拆分]
设输入密文满足
\[
\ct_P \rep P(X^k).
\]
则可通过以下密文操作得到两个分支：

\begin{enumerate}[label=(\arabic*)]
    \item 先做自同构
    \[
    \ct_{-}:=\Auto{1+2n}(\ct_P),
    \]
    其消息为 \(P(-X^k)\)；
    
    \item 再做
    \[
    \cte:=\MulPt{\frac12}\!\bigl(\Add(\ct_P,\ct_{-})\bigr),
    \]
    得到
    \[
    \cte \rep P_{\mathrm{even}}(X^{2k});
    \]

    \item 再做
    \[
    \ctoy:=\MulPt{\frac12}\!\bigl(\Sub(\ct_P,\ct_{-})\bigr),
    \]
    得到
    \[
    \ctoy \rep X^k\,P_{\mathrm{odd}}(X^{2k});
    \]

    \item 若需要去掉前面的 \(X^k\)，则继续做
    \[
    \cto:=\MulPt{-X^{N-k}}(\ctoy),
    \]
    得到
    \[
    \cto \rep P_{\mathrm{odd}}(X^{2k}).
    \]
\end{enumerate}
\end{theorem}

\section{反向方向：由两个分支密文恢复原始密文}

本节给出与上一节完全对应的逆过程。

\subsection{情形 A：奇分支保留 \texorpdfstring{\(Y\)}{Y} 因子}

假设手中已有两个密文：
\[
\cte \rep P_{\mathrm{even}}(X^{2k}),
\qquad
\ctoy \rep X^k\,P_{\mathrm{odd}}(X^{2k}).
\]

定义
\begin{equation}
\ct_P^{(A)}:=\Add(\cte,\ctoy). \label{eq:merge-case-A}
\end{equation}

\begin{theorem}[情形 A 的正确性]
由 \eqref{eq:merge-case-A} 得到的密文满足
\[
\ct_P^{(A)} \rep P(X^k).
\]
\end{theorem}

\begin{proof}
由加法的消息语义，
\[
\ct_P^{(A)}
\rep
P_{\mathrm{even}}(X^{2k})+X^k\,P_{\mathrm{odd}}(X^{2k}).
\]
而根据 \eqref{eq:fft-split}，把其中的 \(Y\) 替换成 \(X^k\) 后即得
\[
P(X^k)=P_{\mathrm{even}}(X^{2k})+X^k\,P_{\mathrm{odd}}(X^{2k}).
\]
所以
\[
\ct_P^{(A)} \rep P(X^k).
\]
\end{proof}

\subsection{情形 B：奇分支已去掉 \texorpdfstring{\(Y\)}{Y} 因子}

假设手中已有
\[
\cte \rep P_{\mathrm{even}}(X^{2k}),
\qquad
\cto \rep P_{\mathrm{odd}}(X^{2k}).
\]

为了恢复 \(P(X^k)\)，需要先把奇分支乘回一个 \(X^k\)。定义
\begin{equation}
\ct_P^{(B)}:=\Add\!\Bigl(\cte,\MulPt{X^k}(\cto)\Bigr). \label{eq:merge-case-B}
\end{equation}

\begin{theorem}[情形 B 的正确性]
由 \eqref{eq:merge-case-B} 得到的密文满足
\[
\ct_P^{(B)} \rep P(X^k).
\]
\end{theorem}

\begin{proof}
先由明文--密文乘法的语义得到
\[
\MulPt{X^k}(\cto)\rep X^k\,P_{\mathrm{odd}}(X^{2k}).
\]
再与 \(\cte\) 相加，可得
\[
\ct_P^{(B)}
\rep
P_{\mathrm{even}}(X^{2k})+X^k\,P_{\mathrm{odd}}(X^{2k}).
\]
由 \eqref{eq:fft-split} 在 \(Y=X^k\) 下的展开，有
\[
P(X^k)=P_{\mathrm{even}}(X^{2k})+X^k\,P_{\mathrm{odd}}(X^{2k}).
\]
因此
\[
\ct_P^{(B)} \rep P(X^k).
\]
\end{proof}

\subsection{反向方向的最终结论}

\begin{theorem}[反向重构]
设 \(\cte,\ctoy,\cto\) 分别满足
\[
\cte \rep P_{\mathrm{even}}(X^{2k}),
\qquad
\ctoy \rep X^k\,P_{\mathrm{odd}}(X^{2k}),
\qquad
\cto \rep P_{\mathrm{odd}}(X^{2k}).
\]
则有两种重构方式：

\begin{enumerate}[label=(\arabic*)]
    \item 若奇分支保留 \(X^k\)，则
    \[
    \ct_P^{(A)}:=\Add(\cte,\ctoy)
    \]
    满足
    \[
    \ct_P^{(A)} \rep P(X^k).
    \]

    \item 若奇分支已去掉 \(X^k\)，则
    \[
    \ct_P^{(B)}:=\Add\!\Bigl(\cte,\MulPt{X^k}(\cto)\Bigr)
    \]
    满足
    \[
    \ct_P^{(B)} \rep P(X^k).
    \]
\end{enumerate}
\end{theorem}

\section{与 FFT 系数向量的对应关系}

虽然本文全程采用的是\textbf{多项式密文视角}，但它与 FFT 递归中熟悉的“偶奇系数向量”完全一致。

定义
\[
\mathbf{e}=(p_0,p_2,\dots,p_{2n-2})^\top\in\RR^n,
\qquad
\mathbf{o}=(p_1,p_3,\dots,p_{2n-1})^\top\in\RR^n.
\]
则
\[
P_{\mathrm{even}}(T)=\sum_{i=0}^{n-1}\mathbf{e}_i\,T^i,
\qquad
P_{\mathrm{odd}}(T)=\sum_{i=0}^{n-1}\mathbf{o}_i\,T^i.
\]

因此，本文正向方向得到的两个密文，本质上正是加密了 FFT 递归中对应于 \(\mathbf e\) 和 \(\mathbf o\) 的两个多项式：
\[
P_{\mathrm{even}}(T)
\quad\text{与}\quad
P_{\mathrm{odd}}(T).
\]

区别仅在于：这里它们以
\[
P_{\mathrm{even}}(X^{2k}),\qquad P_{\mathrm{odd}}(X^{2k})
\]
的形式直接存放在 CKKS 的大环里，而不是先变成槽向量。

\section{结论}

本文给出了一个\textbf{完整、自洽、能直接在 CKKS 密文上执行}的 FFT 式奇偶拆分与重构方案。其核心不是把系数向量显式搬到槽里，再做矩阵抽取；而是直接在多项式环中使用以下精确恒等式：
\[
P(Y)=P_{\mathrm{even}}(Y^2)+Y\,P_{\mathrm{odd}}(Y^2),
\]
\[
P(-Y)=P_{\mathrm{even}}(Y^2)-Y\,P_{\mathrm{odd}}(Y^2).
\]

在固定嵌入
\[
Y=X^{N/(2n)}
\]
的前提下，可以通过一个具体的自同构
\[
\sigma_{1+2n}: X\mapsto X^{1+2n}
\]
实现
\[
Y\mapsto -Y,
\]
然后再通过加法、减法与明文--密文乘法直接得到偶分支与奇分支。反方向则只需做一次加法，或者做一次乘回 \(X^k\) 后再加法，即可恢复原始密文。

因此，从“是否能在 CKKS 密文上严格算通”这一点来看，答案是肯定的，而且整条链条只依赖 CKKS 的标准操作：
\[
\text{自同构} + \text{加法/减法} + \text{明文--密文乘法}.
\]

\end{document}